\section{文档目的与范围}

本文说明「寻迹」校园失物招领平台在开发过程中采用的自动化与协作化手段，以及这些手段
如何作用于需求变更、编码、测试、集成和交付。软件工程化的目的不是增加额外流程，而是
把团队共同遵守的规则固化到仓库、工具和流水线中，使不同成员能够按一致方式开发，并让
重复检查由工具执行。

\begin{table}[H]
  \centering
  \caption{软件工程化手段总览}
  \begin{tabularx}{\textwidth}{L{2.8cm}L{4.1cm}Y}
    \toprule
    方向 & 采用的手段 & 在项目中的作用 \\
    \midrule
    过程组织 & 轻量迭代、需求优先级、任务关闭准则 & 将大型目标拆成可评审、可联调的增量 \\
    团队协作 & Git 分支、Conventional Commit、Pull Request & 隔离并行修改，保留清晰的变更语义 \\
    设计协同 & 需求、API、枚举、数据库文档同步更新 & 保持前端、后端、AI 服务对同一契约的理解一致 \\
    自动检查 & Ruff、Mypy、vue-tsc、Pytest、Vitest & 自动发现格式、类型和业务逻辑问题 \\
    持续集成 & GitHub Actions 五类并行任务 & 在 push 和 Pull Request 时统一执行质量检查 \\
    可重复构建 & uv/pnpm 锁文件、Docker、Compose & 固定依赖解析方式并统一运行环境 \\
    数据演进 & SQLAlchemy 模型与 Alembic 迁移 & 让数据库结构随代码按版本演进 \\
    \bottomrule
  \end{tabularx}
\end{table}

\section{迭代过程与任务协同}

\subsection{轻量迭代方式}

项目采用短周期增量开发。需求按 P0、P1、P2 排序：P0 对应登录、发布、搜索、匹配、
认领、交接和审核等主流程；P1 处理信用、举报、统计和内容治理；P2 用于不影响主流程的
体验增强。每轮迭代从可演示的用户场景出发，同时确定接口、数据、页面和测试任务，避免
不同模块各自推进后再集中拼接。

\begin{longtable}{L{2.0cm}L{3.2cm}L{5.0cm}L{4.1cm}}
  \caption{迭代阶段与协作活动}\label{tab:engineering-iterations}\\
  \toprule
  阶段 & 主要目标 & 团队活动 & 自动化配合 \\
  \midrule
  \endfirsthead
  \multicolumn{4}{c}{\tablename\ \thetable（续）}\\
  \toprule
  阶段 & 主要目标 & 团队活动 & 自动化配合 \\
  \midrule
  \endhead
  \bottomrule
  \endlastfoot
  需求整理 & 明确角色、场景和优先级 & 统一术语、状态和验收条件 & 文档格式与链接检查 \\
  设计协同 & 确定模块、接口和数据关系 & 评审 API、枚举、状态机和迁移方案 & 类型检查与迁移逻辑测试 \\
  功能开发 & 实现可联调的纵向功能 & 分支开发、小步提交、交叉 review & 单元测试、静态检查 \\
  集成联调 & 串联双前端、主后端和 AI 服务 & 按用户场景联调并登记问题 & 前端构建、Compose 配置检查 \\
  交付整理 & 固化版本、说明和演示环境 & 更新文档、脚本与部署配置 & CI 全量检查、PDF 自动构建 \\
\end{longtable}

\subsection{任务关闭准则}

一个开发任务进入合并流程前，需要同时满足以下准则：实现与需求编号或问题描述一致；
接口、枚举、数据库字段和工作流变化同步修改文档；代码通过对应模块的格式、类型和测试
命令；数据库变化附带 Alembic 迁移；界面变化说明交互影响；Pull Request 聚焦一个主题并
由其他成员复核。该准则把「写完代码」转换为团队可共同判断的结束条件。

团队按周协调任务。周初确定本轮用户场景和接口依赖，周中同步跨模块变化及联调问题，
周末按场景走查增量。问题记录包含触发条件、预期行为和实际表现，使负责成员可以直接
定位，而不依赖口头转述。

\section{仓库组织与责任边界}

项目采用 monorepo 管理用户端、管理端、主后端、AI 服务、部署文件和文档。相关变更可以
在同一个 Pull Request 中评审，例如新增一个状态时，可同时查看后端 Schema、前端共享
类型、数据库迁移和接口说明，减少跨仓库版本不同步。

\begin{table}[H]
  \centering
  \caption{模块责任与协作边界}
  \begin{tabularx}{\textwidth}{L{3.0cm}L{4.4cm}Y}
    \toprule
    模块 & 主要责任 & 协作约束 \\
    \midrule
    用户端 & 用户登录、物品发布、搜索、认领与交接交互 & 通过统一 API 层访问后端，共用接口类型与枚举 \\
    管理端 & 内容审核、举报治理、用户管理和统计 & 管理员权限由后端判定，页面不复制业务规则 \\
    主后端 & HTTP 接口、业务状态、事务与数据持久化 & 遵循 router、service、repository 分层 \\
    AI 服务 & 分类、匹配和敏感内容辅助判断 & 通过受鉴权 HTTP 接口调用，不直接访问业务数据库 \\
    部署模块 & 环境模板、镜像构建和六服务编排 & 配置从环境注入，不把本地秘密写入镜像 \\
    文档模块 & 需求、设计、API、测试和报告 & 契约变化与实现置于同一变更中评审 \\
    \bottomrule
  \end{tabularx}
\end{table}

主后端的分层规则是工程约束的重要部分。Router 负责 HTTP 参数、Schema 和身份依赖；
Service 负责业务规则、状态迁移和事务边界；Repository 负责 SQL 查询与持久化。ORM 模型
和 Pydantic Schema 分开维护，使数据库结构、传输格式和业务逻辑各自具有明确职责。

\section{文档与接口契约协作}

\subsection{统一约定}

团队把跨模块信息视为共同契约。JSON 字段使用 camelCase，数据库列使用 snake\_case，
标识符使用 ULID；响应包、分页、时间格式和错误码采用统一约定。角色、物品状态、认领
状态和审核状态使用权威枚举，前端展示值不能自行创造新的字面量。

\begin{table}[H]
  \centering
  \caption{契约变更的同步范围}
  \begin{tabularx}{\textwidth}{L{3.0cm}Y Y}
    \toprule
    变更类型 & 需要同步的对象 & 复核重点 \\
    \midrule
    API 请求或响应 & API 说明、Pydantic Schema、前端 TypeScript 类型 & 字段命名、必填性和错误语义 \\
    枚举或状态机 & 枚举说明、Service 守卫、页面标签与操作入口 & 合法迁移和角色权限 \\
    数据库结构 & ORM、Alembic 迁移、数据库设计说明 & 默认值、索引、唯一约束与历史数据处理 \\
    业务规则 & 需求条目、Service、测试用例和交互说明 & 正常分支、拒绝条件和副作用 \\
    页面流程 & 页面清单、路由、API 调用和验收场景 & 前后端状态是否一致 \\
    \bottomrule
  \end{tabularx}
\end{table}

\subsection{变更影响链}

对跨端需求，团队按「需求场景 $\rightarrow$ 页面操作 $\rightarrow$ API Schema
$\rightarrow$ Service 状态迁移 $\rightarrow$ 数据约束 $\rightarrow$ 自动化测试」的顺序
分析影响。例如调整交接确认规则时，需要同时检查认领详情页、确认接口、认领状态机、
积分副作用和通知逻辑，而不是只修改一个按钮或一个接口。

\section{Git 协作规范}

\subsection{分支与提交}

\term{main} 用于稳定交付，\term{develop} 用于团队集成，功能开发和问题修复分别使用
\term{feature/<name>} 与 \term{fix/<name>}。成员在独立分支中小步提交，减少多人直接修改
同一区域造成的冲突。提交信息采用 Conventional Commit，使历史能够按功能、修复、测试、
文档和工程配置快速筛选。

\begin{verbatim}
feat(FR-02): add lost item publish api
fix(FR-05): correct claim review transition
test(match): cover fallback calculation
docs: update database model
chore(ci): adjust automated checks
\end{verbatim}

\subsection{Pull Request 协作}

Pull Request 关联需求或问题，说明变更范围、受影响模块和本地检查方式。API、枚举、数据库
或工作流变化必须同时包含对应说明；界面变化附关键页面截图。评审者重点检查业务边界、
权限、事务、错误处理和契约一致性。提出修改后，作者以新增提交回应，保留讨论上下文。

仓库通过 \path{.gitattributes} 统一文本行为：一般文本使用 LF，Windows 命令脚本保持
CRLF，图片、PDF、字体和压缩包标记为二进制。\path{.gitignore} 排除环境秘密、依赖目录、
缓存、覆盖率和构建中间产物，使评审内容集中在源代码与必要交付物上。

\section{编码规范与自动检查}

\subsection{Python 模块}

\begin{longtable}{L{3.0cm}L{5.6cm}L{5.6cm}}
  \caption{Python 工程规则与工具}\label{tab:python-engineering}\\
  \toprule
  主题 & 项目约定 & 自动化手段 \\
  \midrule
  \endfirsthead
  \multicolumn{3}{c}{\tablename\ \thetable（续）}\\
  \toprule
  主题 & 项目约定 & 自动化手段 \\
  \midrule
  \endhead
  \bottomrule
  \endlastfoot
  格式与风格 & Python 3.12、100 列、统一 import 与命名 & Ruff lint 与 format check \\
  类型 & 公共函数和返回值使用类型注解，Schema 明确字段约束 & Mypy strict、Pydantic plugin \\
  异步 I/O & 数据库和 HTTP 调用使用显式异步接口 & Ruff ASYNC、Mypy、异步测试 \\
  分层依赖 & Router 不直接访问 Repository，模型与 DTO 分离 & 包结构、代码评审与模块测试 \\
  数据一致性 & 多表写入由 Service 管理事务，关键并发使用锁与唯一约束 & Service 测试、数据库迁移 \\
\end{longtable}

\subsection{前端模块}

双前端使用 Vue 3 Composition API、TypeScript strict、Pinia 和 Vue Router。共享模块保存
接口类型、枚举、标签与格式化函数，两个应用引用同一份源码。统一 Axios 封装负责 API
基址、Bearer token、响应解包、业务错误和未授权会话清理；页面组件只处理展示与交互，
权限和最终业务判断仍由后端负责。

\term{vue-tsc --noEmit} 检查模板和 TypeScript 类型，Vitest 执行工具函数、API 封装和
关键交互逻辑测试，Vite production build 检查生产构建。三类检查同时进入前端 CI 任务，
使开发机与远端流水线采用相同命令。

\section{依赖与构建自动化}

Python 模块使用 uv 管理 \term{pyproject.toml} 和 \term{uv.lock}，两个前端使用 pnpm 管理
\term{package.json} 和各自的 \term{pnpm-lock.yaml}。声明文件表达依赖范围，锁文件固定
实际解析版本；CI 采用冻结安装，避免自动更新依赖导致不同成员或不同日期的构建结果漂移。

\begin{table}[H]
  \centering
  \caption{模块构建约定}
  \begin{tabularx}{\textwidth}{L{3.0cm}L{3.3cm}L{3.9cm}Y}
    \toprule
    模块 & 运行环境 & 依赖工具 & 构建或启动方式 \\
    \midrule
    主后端 & Python 3.12 & uv + uv.lock & uvicorn 启动，Docker 构建镜像 \\
    AI 服务 & Python 3.12 & uv + uv.lock & 独立 uvicorn 进程，Docker 构建镜像 \\
    用户端 & Node 20、pnpm 9 & pnpm lockfile & Vite 构建，Nginx 提供静态文件 \\
    管理端 & Node 20、pnpm 9 & pnpm lockfile & Vite 构建，Nginx 提供静态文件 \\
    完整环境 & Docker/Podman Compose & 镜像 tag 与环境模板 & 编排六个服务并创建持久卷 \\
    \bottomrule
  \end{tabularx}
\end{table}

\section{测试自动化}

测试按模块职责分层。主后端使用 Pytest、pytest-asyncio 和 httpx 的 ASGI transport，
覆盖 Service 规则、状态变化和接口响应；AI 服务使用 Pytest 与 HTTP mock，检查输入
Schema、规则计算、服务鉴权和外部调用处理；双前端使用 Vitest 检查共享工具、API 层和
关键交互。测试所需的数据库会话、对象存储客户端和外部 HTTP 调用由 fixture 或 mock
隔离，使日常测试不依赖手工搭建共享环境。

\begin{table}[H]
  \centering
  \caption{自动化检查的职责分工}
  \begin{tabularx}{\textwidth}{L{3.1cm}L{4.0cm}Y}
    \toprule
    检查层次 & 工具 & 主要发现的问题 \\
    \midrule
    代码风格 & Ruff format、Ruff lint & 格式差异、错误导入、常见 Python 问题 \\
    静态类型 & Mypy、vue-tsc & 参数与返回类型、模板绑定、共享接口不一致 \\
    模块测试 & Pytest、Vitest & 业务规则、状态迁移、工具函数和错误处理 \\
    构建检查 & Vite build、Docker build & 生产包编译和镜像构建问题 \\
    部署检查 & Docker Compose config & 环境变量展开、服务定义和编排语法问题 \\
    \bottomrule
  \end{tabularx}
\end{table}

\section{持续集成流水线}

\subsection{触发规则与执行结构}

GitHub Actions 在 \term{main}、\term{develop} 的 push 和 Pull Request 上触发。流水线按
Backend、AI Service、User App、Admin Web 和 Docker Build 五类任务并行执行；同一分支
有新提交时取消旧运行，避免过时任务继续占用资源。任何任务失败都会使本次检查失败，
所有任务成功后才形成可合并的统一判断。

\begin{figure}[H]
  \centering
  \resizebox{\textwidth}{!}{%
  \begin{tikzpicture}[node distance=1.15cm and 1.0cm]
    \node[state] (event) {push / Pull Request\\main 或 develop};
    \node[state, right=1.2cm of event] (checkout) {检出代码\\安装锁定依赖};

    \node[component, right=2.0cm of checkout, yshift=2.8cm] (backend) {Backend\\Ruff / Mypy / Pytest};
    \node[component, right=2.0cm of checkout, yshift=1.4cm] (ai) {AI Service\\Ruff / Mypy / Pytest};
    \node[component, right=2.0cm of checkout] (userjob) {User App\\类型 / 测试 / 构建};
    \node[component, right=2.0cm of checkout, yshift=-1.4cm] (admin) {Admin Web\\类型 / 测试 / 构建};
    \node[component, right=2.0cm of checkout, yshift=-2.8cm] (docker) {Docker\\Compose 配置 / 镜像};

    \node[decision, right=2.4cm of userjob] (gate) {所有任务\\成功?};
    \node[state, right=1.35cm of gate, yshift=0.85cm] (pass) {允许进入\\合并评审};
    \node[state, right=1.35cm of gate, yshift=-0.85cm] (fail) {返回对应任务\\修改后重跑};

    \draw[flow] (event) -- (checkout);

    \coordinate (fork) at ($(checkout.east)+(0.85cm,0)$);
    \draw[flow] (checkout.east) -- (fork);
    \draw[draw=reportgray,thick] (fork |- backend.west) -- (fork |- docker.west);
    \draw[flow] (fork |- backend.west) -- (backend.west);
    \draw[flow] (fork |- ai.west) -- (ai.west);
    \draw[flow] (fork |- userjob.west) -- (userjob.west);
    \draw[flow] (fork |- admin.west) -- (admin.west);
    \draw[flow] (fork |- docker.west) -- (docker.west);

    \coordinate (gather) at ($(userjob.east)+(0.95cm,0)$);
    \draw[flow] (backend.east) -- (gather |- backend.east);
    \draw[flow] (ai.east) -- (gather |- ai.east);
    \draw[flow] (userjob.east) -- (gather |- userjob.east);
    \draw[flow] (admin.east) -- (gather |- admin.east);
    \draw[flow] (docker.east) -- (gather |- docker.east);
    \draw[draw=reportgray,thick] (gather |- backend.east) -- (gather |- docker.east);
    \draw[flow] (gather) -- (gate.west);

    \draw[flow] (gate) -- node[above,sloped]{是} (pass);
    \draw[flow] (gate) -- node[below,sloped]{否} (fail);
  \end{tikzpicture}}
  \caption{持续集成流水线}
  \label{fig:ci-pipeline}
\end{figure}

\subsection{各任务采用的自动化步骤}

\begin{longtable}{L{2.8cm}L{7.2cm}L{4.6cm}}
  \caption{CI 任务与检查内容}\label{tab:ci-jobs-engineering}\\
  \toprule
  任务 & 自动化步骤 & 输出用途 \\
  \midrule
  \endfirsthead
  \multicolumn{3}{c}{\tablename\ \thetable（续）}\\
  \toprule
  任务 & 自动化步骤 & 输出用途 \\
  \midrule
  \endhead
  \bottomrule
  \endlastfoot
  Backend & uv 冻结安装；Ruff lint/format；Mypy；Pytest coverage & 统一 Python 规范、类型、业务测试和覆盖率记录 \\
  AI Service & uv 安装；Ruff lint/format；Mypy；Pytest & 检查独立服务的代码、类型和接口行为 \\
  User App & pnpm 冻结安装；vue-tsc；Vitest；Vite build & 检查用户端类型、逻辑和生产包 \\
  Admin Web & pnpm 冻结安装；vue-tsc；Vitest；Vite build & 检查管理端类型、逻辑和生产包 \\
  Docker Build & Compose 配置展开；Backend 与 AI 镜像构建 & 检查部署配置和镜像构建链 \\
\end{longtable}

Backend 任务生成覆盖率 XML 作为流水线产物，便于查看未覆盖模块；Compose 检查使用专用
占位配置验证必填环境变量和编排语法，不使用开发者本地的环境文件。流水线只负责提供
一致判断，具体修改仍由对应模块成员在原分支完成。

\section{数据库与环境自动化}

\subsection{数据库迁移}

数据库结构通过 SQLAlchemy 模型和 Alembic 迁移共同管理。开发者修改模型后生成迁移
初稿，再人工检查字段类型、默认值、索引、唯一约束和历史数据转换。部署时先执行
\term{alembic upgrade head}，再启动主后端；迁移历史保持单一 head，使不同环境能够从
明确版本按顺序升级。

涉及数据结构的 Pull Request 同时包含 ORM、迁移和设计说明。自动生成只负责比较结构，
评审者负责判断数据语义。例如建立唯一约束前，需要先处理历史重复记录；修改状态字段时，
需要给出旧值到新值的转换规则。

\subsection{环境与容器}

运行配置由环境变量注入，仓库保留不含真实秘密的示例模板。本地完整栈由 Compose 编排
MySQL、MinIO、AI Service、Backend、User Web 和 Admin Web；后端容器在启动应用前执行
迁移，前端容器使用 Nginx 提供静态资源并代理 API。镜像构建和本地启动采用同一依赖
锁文件，减少「开发机可以运行、交付环境不能运行」的差异。

\section{典型变更的工程化流程}

\begin{longtable}{L{3.0cm}L{8.2cm}L{3.4cm}}
  \caption{典型变更处理流程}\label{tab:change-flow}\\
  \toprule
  变更类型 & 处理步骤 & 自动检查 \\
  \midrule
  \endfirsthead
  \multicolumn{3}{c}{\tablename\ \thetable（续）}\\
  \toprule
  变更类型 & 处理步骤 & 自动检查 \\
  \midrule
  \endhead
  \bottomrule
  \endlastfoot
  新增 API & 确认需求与错误语义；编写 Schema、Service、Router；同步前端类型和接口说明；补充测试 & Ruff、Mypy、Pytest、vue-tsc、Vitest \\
  修改状态机 & 更新合法迁移；检查权限和副作用；同步枚举与页面操作；覆盖拒绝分支 & 后端测试、前端类型和构建 \\
  修改数据库 & 更新 ORM；编写 Alembic；说明数据转换；检查单一 head；同步设计文档 & 迁移检查、Pytest、Docker build \\
  调整页面 & 确认用户场景；修改组件、Store 与 API 调用；检查加载、空态和错误反馈 & vue-tsc、Vitest、Vite build \\
  调整部署 & 修改环境模板或 Compose；说明变量含义、默认值和服务依赖 & Compose config、Docker build \\
\end{longtable}

当自动检查失败时，团队按任务名称定位到模块，再根据日志修正实现或测试。修正仍提交到
原分支，由流水线重新执行完整任务。这个反馈过程将代码评审和机器检查组合起来：工具处理
重复且可判定的规则，成员集中评审需求理解、架构边界和业务语义。

\section{结论}

本文说明了寻迹项目实际采用的软件工程化方法。自动化方面，项目使用依赖锁、静态检查、
类型检查、模块测试、生产构建、容器构建、Compose 配置检查和 Alembic 迁移，并通过
GitHub Actions 统一执行；协作化方面，项目使用轻量迭代、任务关闭准则、模块责任边界、
分支与提交规范、Pull Request 评审以及文档和代码同步机制。

这些方法覆盖从需求进入开发到代码合并和环境构建的全过程，使团队能够并行工作，并以
一致的规则处理接口、数据、测试和交付变化。
